\renewcommand{\Anno}{2013}

\renewcommand{\NumeroProblema}{4}

\renewcommand{\Testo}{
Gli studenti della Scuola Sant'Anna sono riusciti a produrre tre ottime birre, tutte dello stesso gusto ma con gradazioni alcoliche diverse: una completamente analcolica, una con gradazione alcolica di $5^{\circ}$, una con gradazione alcolica di $10^{\circ}$.  Al bar della Scuola l'ordinazione standard è: ``un litro di birra con gradazione di $g$ gradi''.  Le gradazioni richieste dai clienti sono numeri reali $g$ distribuiti nell'intervallo $[0,10]$ in maniera completamente casuale.  Il barman, miscelando i tre tipi di birra, cerca di soddisfare la richiesta.

Ad un certo istante, nel bar restano 1 litro di birra a $0^{\circ}$, 1 litro di birra a $5^{\circ}$ ed 1 litro di birra a $10^{\circ}$. Entra un cliente e chiede 1 litro di birra a $2.5^{\circ}$.
\begin{enumerate}
	\item[(a)] Determinare in che modo il barman dovrà miscelare le tre birre per avere la \emph{massima} probabilità di poter soddisfare anche il cliente successivo.  Determinare anche, posto che il barman opti per tale strategia ottimale, per quali valori di $g$ il cliente successivo potrà essere accontentato.

	\item[(b)] Viceversa, determinare in che modo il barman dovrà miscelare le tre birre per avere la \emph{minima} probabilità di poter soddisfare anche il cliente successivo, precisando anche in questa ipotesi pessimistica per quali valori di $g$ il cliente successivo potrà essere accontentato.
\end{enumerate}
}

\renewcommand{\Risposta}{
\begin{enumerate}
	\renewcommand{\labelenumi}{(\alph{enumi})}
	\item[(a)] Per avere la massima probabilità di accontentare anche il cliente successivo, il barman deve miscelare mezzo litro di birra analcolica e mezzo litro di birra a $5^{\circ}$.  In tal caso potrà accontentare il cliente successivo se e solo se $2.5\leq g\leq 10$, cioè nel 75\% dei casi.

	\item[(b)] Per avere la minima probabilità di accontentare anche il cliente successivo, il barman deve miscelare 3/4 di litro di birra analcolica e 1/4 di litro di birra a $10^{\circ}$.  In tal caso potrà accontentare il cliente successivo se e solo se $3.75\leq g\leq 8.75$, cioè nel 50\% dei casi.
\end{enumerate}
}

\renewcommand{\Aiutino}{
Supponiamo che il barman ottenga la prima birra miscelando $a$ litri di quella analcolica, $b$ litri di quella a $5^{\circ}$, e $c$ litri di quella a $10^{\circ}$.  Quali relazioni devono verificare $a$, $b$, $c$?

Si tratta poi di dimostrare che il cliente successivo potrà essere accontentato se e solo se richiederà una gradazione $g$ con $a/2\leq g\leq 1-c/2$.  Non resta ora che determinare per quali valori dei parametri tale intervallo ha ampiezza massima o minima.
}

\renewcommand{\SoluzioneUfficiale}{\TitoloSoluzione
Per comodità dividiamo tutte le gradazioni per 10: così le 3 birre standard hanno gradazione 0, 1/2, 1, mentre il primo cliente richiede gradazione 1/4. Supponiamo che il barman accontenti il primo cliente usando $a$ litri della birra a gradazione 0, $b$ litri della birra a gradazione 1/2, e $c$ litri della birra a gradazione 1. Chiaramente $a\geq 0$, $b\geq 0$, $c\geq 0$; inoltre
\begin{equation}
	a+b+c=1
	\label{birre:1}
\end{equation}
perché la somma deve fare 1 litro di birra, e
\begin{equation}
	\frac{b}{2}+c=a\cdot 0+b\cdot\frac{1}{2}+c\cdot 1=\frac{1}{4},
	\label{birre:1/4}
\end{equation}
in quanto la gradazione totale deve essere 1/4.

Mostriamo che il cliente successivo potrà essere accontentato se e solo se richiederà una gradazione $g$ con
\begin{equation}
	\frac{a}{2}\leq g\leq 1-\frac{c}{2}.
	\label{birre:g}
\end{equation}

Se dimostriamo questo avremo che le gradazioni disponibili per il cliente successivo sono un intervallo di ampiezza
$$1-\frac{c}{2}-\frac{a}{2}=1-\frac{a+c}{2}=\frac{1+b}{2},$$
in cui l'ultima uguaglianza segue dalla (\ref{birre:1}). Ora dalla
(\ref{birre:1/4}) deduciamo immediatamente che $0\leq b\leq 1/2$. A
questo punto è facile rispondere alle domande poste dal problema come segue.
\begin{itemize}
	\item  Il cliente successivo ha a disposizione un intervallo di gradazioni di ampiezza massima quando $b$ è il massimo possibile, cioè quando $b=1/2$. Tenendo conto della
	(\ref{birre:1}) e della (\ref{birre:1/4}), questo corrisponde ad
	avere $a=1/2$, $b=1/2$, $c=0$, cioè al caso in cui il barman ottiene la birra per il primo cliente mescolando mezzo litro di quella analcolica e mezzo litro di quella a gradazione 1/2.
	Tenendo conto della (\ref{birre:g}), in tal caso il cliente
	successivo sarà accontentabile se e solo se $1/4\leq g\leq
	1$, cioè nel 75\% dei casi.

	\item  Il cliente successivo ha a disposizione un intervallo di gradazioni di ampiezza minima quando $b$ è il minimo possibile, cioè quando $b=0$. Tenendo conto della
	(\ref{birre:1}) e della (\ref{birre:1/4}), questo corrisponde ad
	avere $a=3/4$, $b=0$, $c=1/4$, cioè al caso in cui il barman ottiene la birra per il primo cliente mescolando 3/4 di litro di quella analcolica e 1/4 di litro di quella a gradazione 1.
	Tenendo conto della (\ref{birre:g}), in tal caso il cliente
	successivo sarà accontentabile se e solo se $3/8\leq g\leq
	7/8$, cioè nel 50\% dei casi.
\end{itemize}

Resta dunque da dimostrare la (\ref{birre:g}).  Dopo aver accontentato
il primo cliente, restano $(1-a)$ litri a gradazione 0, $(1-b)$ litri a gradazione 1/2, e $(1-c)$ litri a gradazione 1.  La birra da 1 litro più leggera che si può produrre in queste condizioni è quella ottenuta usando tutta la birra a gradazione 0 disponibile, cioè $(1-a)$ litri, e la restante parte, cioè $a$ litri, a gradazione 1/2.  Questa scelta è possibile in quanto sono rimasti almeno $a$ litri di birra a gradazione 1/2, dal momento che $a\leq 1-b$
per via della (\ref{birre:1}).  La birra ottenuta in questo modo ha
gradazione
$$(1-a)\cdot 0+a\cdot \frac{1}{2}=\frac{a}{2},$$
il che giustifica intanto la disuguaglianza $g\geq a/2$, e mostra come produrre 1 litro di birra a gradazione esattamente $a/2$.

Simmetricamente, la birra da 1 litro più alcolica che il cliente successivo può avere è quella ottenuta usando tutta la birra a gradazione 1 disponibile, cioè $(1-c)$ litri, e la restante parte, cioè $c$ litri, a gradazione 1/2.  Questa scelta è possibile in quanto $c\leq 1-b$ per lo stesso motivo di prima, e porta ad una birra di gradazione totale
$$c\cdot\frac{1}{2}+(1-c)\cdot 1=1-\frac{c}{2}.$$

Questo giustifica la disuguaglianza $g\leq 1-c/2$, e mostra come produrre 1 litro di birra a gradazione esattamente $1-c/2$.

Resta da dimostrare che tutte le gradazioni intermedie tra $a/2$ e $1-c/2$ sono pure ottenibili.  A tal fine il barman può produrre 1 litro di birra a gradazione $a/2$ procedendo come descritto, ed 1 litro di birra a gradazione $1-c/2$ procedendo come descritto.  Si noti che questo lo porta ad usare $a$ litri di birra a gradazione 1/2 per produrre quella a gradazione $a/2$, e $c$ litri di birra a gradazione 1/2 per produrre quella a gradazione $1-c/2$.  Sommati ai $b$ litri di birra a gradazione 1/2 già utilizzati per il primo cliente, il totale fa esattamente 1.  A questo punto è evidente
che tutte le gradazioni che soddisfano la (\ref{birre:g}) possono
essere ottenute miscelando in maniera opportuna le 2 birre
``estremali'' appena prodotte.  }

\renewcommand{\SoluzioniAlternative}{
}

\renewcommand{\Osservazioni}{
È possibile, seguendo la strategia a due passi descritta nella soluzione, determinare in funzione di $a$, $b$, $c$, $g$ le proporzioni secondo cui miscellare le 3 birre all'arrivo del cliente successivo. Con qualche calcolo, che è comunque istruttivo svolgere, si ottiene che le birre di gradazione 0, 1/2, 1 devono essere mescolate, rispettivamente, nelle seguenti quantità:
$$a'=\frac{(1-a)(2-c-2g)}{1+b},
\hspace{3em}
b'=\frac{2(a-ac-ag+cg)}{1+b},
\hspace{3em}
c'=\frac{(1-c)(2g-a)}{1+b}.$$

Con qualche calcolo si verifica che, nelle ipotesi (\ref{birre:g}), si
ha che $0\leq a'\leq 1-a$, $0\leq b'\leq 1-b$, $0\leq c'\leq 1-c$, ed inoltre $a'+b'+c'=1$ e $b'/2+c'=g$.
}

\renewcommand{\Commenti}{
}
